\documentclass{article}
\usepackage[submission]{colm2026_conference}

\usepackage{microtype}
\usepackage{hyperref}
\usepackage{url}
\usepackage{booktabs}
\usepackage{amsmath}
\usepackage{amssymb}
\usepackage{graphicx}
\usepackage{multirow}
\usepackage{xcolor}
\usepackage{subcaption}
\usepackage{enumitem}
\usepackage{amsthm}
\newtheorem{theorem}{Theorem}

\usepackage{lineno}

\definecolor{darkblue}{rgb}{0, 0, 0.5}
\hypersetup{colorlinks=true, citecolor=darkblue, linkcolor=darkblue, urlcolor=darkblue}

\title{Appendix: How Reliable Are VLM Judges? Conformal Prediction Reveals \\ Task-Dependent Uncertainty in Multimodal Evaluation}

\author{Anonymous authors\\Paper under double-blind review}

\begin{document}

\maketitle

\appendix

\section{Evaluation prompts}
\label{app:prompts}

\subsection{MLLM-as-a-Judge CoT prompt}

We use a chain-of-thought prompt adapted from the MLLM-as-a-Judge dataset~\citep{chen2024mllmasajudge}. The prompt follows the ``COT figure'' setting that forces the judge to analyze the image first, then evaluate the response, then provide a final score:

\begin{quote}
\small
\textit{Please serve as an unbiased judge in assessing the quality of the response from an AI assistant regarding the user's instruction and the provided image.}

\textit{Evaluation Steps: Please examine the provided image attentively. Begin by conducting a comprehensive analysis of the figure provided. Then, utilize the insights from your analysis to critically evaluate the response. Finally, based on your figure analysis and response evaluation, form a well-reasoned judgement.}

\textit{Scoring Rubric:}
\textit{Poor (1): The response significantly deviates from the user's instruction and fails to address the query effectively. It shows a lack of relevance, accuracy, and comprehensiveness. Creativity and granularity are absent or poorly executed.}
\textit{Fair (2): The response addresses the user's instruction partially, with evident shortcomings in relevance, accuracy, or comprehensiveness. It lacks depth in creativity and granularity.}
\textit{Average (3): The response adequately addresses the user's instruction, showing a fair level of relevance, accuracy, and comprehensiveness. It reflects a basic level of creativity and granularity but may lack sophistication.}
\textit{Good (4): The response is well-aligned with the user's instruction, demonstrating a high degree of relevance, accuracy, and comprehensiveness. It shows creativity and nuanced understanding.}
\textit{Excellent (5): The response perfectly adheres to the user's instruction, excelling in relevance, accuracy, comprehensiveness, creativity, and granularity.}

\textit{Notice: In your evaluation, weigh factors such as relevance, accuracy, comprehensiveness, creativity, and the granularity of the response. Do not allow the length of the response to influence your evaluation. Be as objective as possible.}

\textit{Here is the input:}
\textit{[The Start of User Instruction] \{question\} [The End of User Instruction]}
\textit{[The Start of Assistant's Answer] \{actor\_answer\} [The End of Assistant's Answer]}

\textit{First, provide your analysis of the image and the response. Then, at the very end of your response, provide your final score in exactly this format on its own line: ``Score: X'' where X is a single number from 1 to 5.}
\end{quote}

\subsection{Polaris captioning prompt}

For the Polaris captioning dataset, we adapt the prompt structure for caption evaluation:

\begin{quote}
\small
\textit{Please serve as an unbiased judge in assessing the quality of an image caption generated by an AI model.}

\textit{Evaluation Steps: Please examine the provided image attentively. Then evaluate how well the generated caption describes the image content.}

\textit{Scoring Rubric:}
\textit{Poor (1): The caption is completely inaccurate, irrelevant, or fails to describe the image content. Major factual errors or hallucinations are present.}
\textit{Fair (2): The caption partially describes the image but has significant inaccuracies, missing key elements, or includes incorrect details.}
\textit{Average (3): The caption adequately describes the main content of the image with fair accuracy, but may miss some details or include minor inaccuracies.}
\textit{Good (4): The caption accurately describes the image content with good detail and relevance, with only minor omissions.}
\textit{Excellent (5): The caption perfectly describes the image content with high accuracy, comprehensive detail, and excellent relevance to what is shown.}

\textit{Provide your analysis of how well the caption matches the image, then provide your final score in exactly this format on its own line: ``Score: X'' where X is a single number from 1 to 5.}
\end{quote}

\section{Additional per-dataset results}
\label{app:per_dataset}

\subsection{All judges per-dataset width comparison}

Table~\ref{tab:all_judges_per_dataset} presents the full per-dataset R2CCP width comparison across all three judges. Gemini produces narrower intervals on 10 of 14 datasets, with the largest improvements on Vision-Heavy tasks.

\begin{table}[t]
\centering
\caption{Per-dataset R2CCP raw width for all three judges. Bold indicates narrowest width for each dataset. Gemini produces narrower intervals on 10 of 14 datasets, primarily Vision-Heavy tasks.}
\label{tab:all_judges_per_dataset}
\small
\begin{tabular}{lcccc}
\toprule
\textbf{Dataset} & \textbf{LLaVA} & \textbf{Phi-4} & \textbf{Gemini} & \textbf{Cat.} \\
\midrule
AesBench & 2.08 & \textbf{1.98} & 2.14 & Aesth. \\
MM-Vet & \textbf{2.18} & 2.39 & 2.29 & Gen. \\
WIT & \textbf{2.38} & 2.48 & 2.39 & Know. \\
COCO & 2.43 & 2.51 & \textbf{2.40} & Gen. \\
Mind2Web & \textbf{2.69} & 2.74 & 2.92 & Know. \\
Conc. Cap. & 2.70 & \textbf{2.70} & 2.80 & Know. \\
TextVQA & 2.81 & 3.10 & \textbf{2.33} & Vision \\
LLaVA-B. & 2.92 & 3.07 & \textbf{2.91} & Gen. \\
VisitBench & 2.96 & 2.90 & \textbf{2.85} & Gen. \\
ChartQA & 3.08 & 3.52 & \textbf{2.83} & Vision \\
ScienceQA & 3.27 & 3.25 & \textbf{2.77} & Vision \\
MathVista & 3.37 & 3.42 & \textbf{3.39} & Vision \\
DiffusionDB & 3.41 & \textbf{3.22} & 3.38 & Aesth. \\
Infograph. & 3.50 & 3.58 & \textbf{2.93} & Vision \\
\midrule
\textbf{Overall} & 3.05 & 3.13 & \textbf{2.85} & --- \\
\bottomrule
\end{tabular}
\end{table}

\subsection{Per-dataset relaxed accuracy}

Table~\ref{tab:relaxed_accuracy} presents the relaxed accuracy ($\pm 1$) and MAE for each dataset, sorted by $\pm 1$ accuracy. Tasks with higher accuracy generally produce narrower conformal intervals, but the relationship is not monotonic due to the ranking-scoring decoupling.

\begin{table}[t]
\centering
\caption{Per-dataset relaxed accuracy and MAE for LLaVA-Critic-7B. Higher accuracy generally corresponds to narrower intervals but the relationship is not monotonic.}
\label{tab:relaxed_accuracy}
\small
\begin{tabular}{lrcccc}
\toprule
\textbf{Dataset} & \textbf{N} & \textbf{Exact} & $\pm$\textbf{1} & \textbf{MAE} & \textbf{Bias} \\
\midrule
AesBench & 392 & 41.3\% & 88.3\% & 0.732 & $+0.49$ \\
COCO & 397 & 34.5\% & 84.9\% & 0.854 & $+0.04$ \\
Conc. Cap. & 398 & 35.4\% & 81.7\% & 0.872 & $+0.60$ \\
LLaVA-B. & 396 & 31.8\% & 81.1\% & 0.922 & $+0.55$ \\
MM-Vet & 258 & 37.2\% & 80.2\% & 0.926 & $-0.01$ \\
TextVQA & 399 & 36.3\% & 79.2\% & 0.942 & $+0.30$ \\
WIT & 399 & 33.8\% & 78.7\% & 0.912 & $+0.51$ \\
VisitBench & 397 & 31.7\% & 75.1\% & 1.018 & $+0.44$ \\
Mind2Web & 398 & 29.6\% & 74.4\% & 1.008 & $+0.19$ \\
ChartQA & 400 & 30.2\% & 73.2\% & 1.085 & $+0.31$ \\
ScienceQA & 396 & 26.5\% & 69.7\% & 1.200 & $+0.30$ \\
Infograph. & 398 & 25.1\% & 69.3\% & 1.191 & $+0.18$ \\
MathVista & 789 & 32.3\% & 65.1\% & 1.235 & $+0.33$ \\
DiffusionDB & 299 & 23.7\% & 57.2\% & 1.385 & $+1.23$ \\
\bottomrule
\end{tabular}
\end{table}

\subsection{Confusion matrices}

Table~\ref{tab:confusion} shows the row-normalized confusion matrices for all three judges, revealing distinct error patterns: LLaVA-Critic shows broad dispersion, Phi-4 exhibits gravity toward score 4, and Gemini polarizes between scores 1 and 5.

\begin{table}[t]
\centering
\caption{Row-normalized confusion matrices (\%, GT rows, Pred columns). LLaVA shows broad dispersion, Phi-4 gravitates to 4, Gemini polarizes to 1 and 5.}
\label{tab:confusion}
\small
\begin{tabular}{l|ccccc|c}
\toprule
& \textbf{P=1} & \textbf{P=2} & \textbf{P=3} & \textbf{P=4} & \textbf{P=5} & $\pm$\textbf{1} \\
\midrule
\multicolumn{7}{l}{\emph{LLaVA-Critic-7B}} \\
GT=1 & 32.6 & 12.9 & 17.0 & 23.4 & 14.0 & 45.6 \\
GT=2 & 17.0 & 11.7 & 21.9 & 32.2 & 17.1 & 50.6 \\
GT=3 & 5.2 & 7.0 & 18.6 & 42.0 & 27.2 & 67.6 \\
GT=4 & 3.3 & 3.6 & 13.4 & 38.9 & 40.8 & 93.1 \\
GT=5 & 1.9 & 2.2 & 10.7 & 38.4 & 46.7 & 85.2 \\
\midrule
\multicolumn{7}{l}{\emph{Phi-4-reasoning-15B}} \\
GT=1 & 21.6 & 18.1 & 13.8 & 33.6 & 13.0 & 39.6 \\
GT=2 & 12.8 & 18.4 & 21.5 & 37.9 & 9.4 & 52.7 \\
GT=3 & 6.4 & 13.3 & 28.1 & 44.5 & 7.7 & 85.9 \\
GT=4 & 2.4 & 10.2 & 19.5 & 56.2 & 11.8 & 87.4 \\
GT=5 & 2.5 & 4.6 & 11.4 & 56.9 & 24.6 & 81.4 \\
\midrule
\multicolumn{7}{l}{\emph{Gemini 2.5 Flash}} \\
GT=1 & 61.7 & 15.5 & 6.4 & 5.2 & 11.3 & 77.1 \\
GT=2 & 38.7 & 24.7 & 9.3 & 9.9 & 17.5 & 72.7 \\
GT=3 & 17.6 & 21.3 & 12.0 & 18.8 & 30.2 & 52.1 \\
GT=4 & 10.4 & 13.7 & 9.0 & 15.0 & 52.0 & 76.0 \\
GT=5 & 7.9 & 10.3 & 6.7 & 11.4 & 63.7 & 75.0 \\
\bottomrule
\end{tabular}
\end{table}

\section{CP methods across all three judges}
\label{app:all_methods}

Table~\ref{tab:all_methods_all_judges} compares the top conformal methods across all three judges on MLLM-as-a-Judge. R2CCP and CHR consistently emerge as the best methods, with Gemini producing the narrowest intervals for both.

\begin{table}[t]
\centering
\caption{Top CP methods across judges (raw coverage/width, 10 seeds). R2CCP and CHR are consistently the best methods. Gemini produces the narrowest intervals for both methods.}
\label{tab:all_methods_all_judges}
\small
\begin{tabular}{lcccccc}
\toprule
& \multicolumn{2}{c}{\textbf{LLaVA-Critic}} & \multicolumn{2}{c}{\textbf{Phi-4}} & \multicolumn{2}{c}{\textbf{Gemini}} \\
\cmidrule(lr){2-3} \cmidrule(lr){4-5} \cmidrule(lr){6-7}
\textbf{Method} & Cov. & Width & Cov. & Width & Cov. & Width \\
\midrule
R2CCP & .900 & 3.05 & .891 & 3.13 & .898 & \textbf{2.85} \\
CHR & .880 & 2.97 & .896 & 3.16 & .880 & \textbf{2.80} \\
B. LCP & .863 & 3.02 & .875 & 3.24 & .863 & \textbf{2.88} \\
Naive CP & .895 & 3.23 & .897 & 3.30 & .893 & \textbf{3.08} \\
CQR & .996 & 3.95 & .997 & 3.98 & .996 & 3.95 \\
\bottomrule
\end{tabular}
\end{table}

\section{Polaris detailed results}
\label{app:polaris}

\subsection{Polaris confusion matrix}

Table~\ref{tab:polaris_confusion} shows the confusion matrix for LLaVA-Critic-7B on Polaris. The judge achieves near-perfect identification of GT$=$1 captions (99.6\% exact accuracy), demonstrating that it can reliably identify clearly bad captions. The most confused categories are GT$=$2 and GT$=$3, where the judge tends to underscore by assigning GT$=$1.

\begin{table}[t]
\centering
\caption{Confusion matrix for LLaVA-Critic-7B on Polaris captioning dataset (\%). GT=1 is nearly perfect (99.6\%).}
\label{tab:polaris_confusion}
\small
\begin{tabular}{l|ccccc|c}
\toprule
& \textbf{P=1} & \textbf{P=2} & \textbf{P=3} & \textbf{P=4} & \textbf{P=5} & $\pm$\textbf{1} \\
\midrule
GT=1 & 99.6 & 0.1 & 0.2 & 0.1 & 0.0 & 99.7 \\
GT=2 & 74.2 & 12.9 & 5.7 & 4.8 & 2.4 & 92.7 \\
GT=3 & 18.6 & 13.4 & 32.1 & 30.8 & 5.2 & 76.2 \\
GT=4 & 1.6 & 4.0 & 14.6 & 64.2 & 15.6 & 94.4 \\
GT=5 & 0.8 & 2.0 & 5.6 & 63.9 & 27.8 & 91.7 \\
\bottomrule
\end{tabular}
\end{table}

\subsection{Polaris error-bin CP coverage}

Table~\ref{tab:polaris_error_bins} shows CP coverage by error magnitude on Polaris. The tight intervals (avg.\ width 0.68) mean that even $\pm 1$ errors have lower raw coverage (58.1\%) than on MLLM-Judge (98.7\%), because the intervals are so narrow they don't always extend to neighboring scores. After boundary adjustment, coverage recovers to 99.4\%.

\begin{table}[t]
\centering
\caption{CP coverage by error magnitude on Polaris (R2CCP, LLaVA-Critic-7B). Tight intervals mean lower raw coverage for $\pm 1$ errors, but boundary adjustment recovers to 99.4\%.}
\label{tab:polaris_error_bins}
\small
\begin{tabular}{lrcc}
\toprule
\textbf{Error} & \textbf{\% samples} & \textbf{Cov. (raw)} & \textbf{Cov. (adj.)} \\
\midrule
Exact (0) & 80.9\% & 98.2\% & 99.6\% \\
$\pm 1$ & 14.6\% & 58.1\% & 99.4\% \\
$\pm 2$ & 4.1\% & 50.0\% & 80.7\% \\
$\pm 3$ & 0.5\% & 0.6\% & 52.4\% \\
\bottomrule
\end{tabular}
\end{table}

\section{Naive Split CP vs.\ R2CCP per dataset}
\label{app:naive_vs_r2ccp}

Table~\ref{tab:naive_vs_r2ccp} demonstrates that R2CCP consistently outperforms the Naive Split CP baseline across all 14 datasets, with width improvements ranging from 0.15 to 1.21 score points.

\begin{table}[t]
\centering
\caption{Naive Split CP vs.\ R2CCP per dataset (raw width, LLaVA-Critic-7B). R2CCP produces narrower intervals on all 14 datasets.}
\label{tab:naive_vs_r2ccp}
\small
\begin{tabular}{lccc}
\toprule
\textbf{Dataset} & \textbf{Naive} & \textbf{R2CCP} & \textbf{$\Delta$} \\
\midrule
AesBench & 2.53 & \textbf{2.08} & $-0.45$ \\
MM-Vet & 3.39 & \textbf{2.18} & $-1.21$ \\
WIT & 2.92 & \textbf{2.38} & $-0.54$ \\
COCO & 2.86 & \textbf{2.43} & $-0.43$ \\
Mind2Web & 3.13 & \textbf{2.69} & $-0.44$ \\
Conc. Cap. & 3.11 & \textbf{2.70} & $-0.41$ \\
TextVQA & 3.19 & \textbf{2.81} & $-0.38$ \\
LLaVA-B. & 3.28 & \textbf{2.92} & $-0.36$ \\
VisitBench & 3.37 & \textbf{2.96} & $-0.41$ \\
ChartQA & 3.49 & \textbf{3.08} & $-0.41$ \\
ScienceQA & 3.71 & \textbf{3.27} & $-0.44$ \\
MathVista & 3.64 & \textbf{3.37} & $-0.27$ \\
DiffusionDB & 3.56 & \textbf{3.41} & $-0.15$ \\
Infograph. & 3.67 & \textbf{3.50} & $-0.17$ \\
\midrule
\textbf{Average} & 3.27 & \textbf{2.84} & $-0.42$ \\
\bottomrule
\end{tabular}
\end{table}

\section{Effect of chain-of-thought prompting}
\label{app:cot}

Table~\ref{tab:cot_effect} compares the generic prompt (no CoT) vs.\ our CoT prompt for LLaVA-Critic-7B. CoT improves correlation metrics by 6--13\% and---critically for conformal prediction---reduces overconfidence by 50\%, giving CP methods more distributional signal in the logprobs.

\begin{table}[t]
\centering
\caption{Effect of chain-of-thought prompting on LLaVA-Critic-7B performance and logprob calibration. CoT halves the overconfidence rate, providing better signal for conformal prediction.}
\label{tab:cot_effect}
\small
\begin{tabular}{lcc}
\toprule
\textbf{Metric} & \textbf{Generic prompt} & \textbf{CoT prompt} \\
\midrule
Pearson $\rho$ & .380 & \textbf{.402} ($+5.8\%$) \\
Spearman $\rho_s$ & .314 & \textbf{.356} ($+13.4\%$) \\
Kendall $\tau$ & .270 & \textbf{.300} ($+11.1\%$) \\
Exact accuracy & \textbf{33.9\%} & 32.2\% \\
MAE & 1.058 & \textbf{1.031} \\
\midrule
Overconf.\ ($>$0.99) & 51.4\% & \textbf{25.6\%} ($-50\%$) \\
Overconf.\ ($>$0.999) & 33.2\% & \textbf{6.6\%} ($-80\%$) \\
\bottomrule
\end{tabular}
\end{table}

\bibliography{references}
\bibliographystyle{colm2026_conference}

\end{document}
